\section{Related Work}

\textbf{Out-of-Core and Semi-External GPS:} \PM{talk here about how these systems have been optimized for HDD/fast SSDs and need the developer to know in detail *how* to optimize their system}
Starting with GraphChi \cite{kyrola2012graphchi} and X-Stream\cite{roy2013x}, several out-of-core graph processing systems \cite{zhu2015gridgraph, vora2019lumos} have been proposed.
These systems are more resource efficient, easier to manage and more performant than distributed graph processing systems, but often perform worse than single-node shared memory systems such as Galois \cite{pingali2011tao} and Ligra\cite{shun2013ligra}.
GraphChi was designed to keep the IO as sequential as possible and, to achieve that, using it requires preprocessing the data and maintaining shards that minimize random reads or writes during its parallel sliding window execution model.
Similarly, X-Stream accesses the edges from disk sequentially, reading more data than is necessary to avoid the prohibitive cost of random reads from a hard disk.
Advances in storage research are also indicated in the development of out-of-core graph processing systems.
Flashgraph\cite{zheng2015flashgraph} and Graphene\cite{liu2017graphene} are systems that have been designed to leverage the fast random IO supported by SSDs; they strive to minimize write amplification instead of reducing random IO.
Blaze\cite{kim2022blaze} has been similarly designed for fast flash devices; their design goal is to maximize the bandwidth utilization and leverage the fast random IO supported by flash devices. 

Another explicit design goal of these early out-of-core systems was to keep memory demands of graph processing low, and the memory budget these applications require is often an explicit tunable parameter.
Over time, the price of RAM has fallen, but overprovisioning memory remains a leading expense for system operators and is sensitive to the global supply of cheap silicon.
Semi-External systems take an alternative approach: instead of keeping the entire graph in shared memory, they hold the graph vertex data and the algorithmic state in memory.
Several systems, such as CLIP\cite{ai2018clip} and Blaze\cite{kim2022blaze}, use the semi-external processing model, and they perform better than out-of-core systems by minimizing disk I/O for managing vertex data.
The choice between these two paradigms depends on the memory budget of the system and the size of the graph being analyzed.
\project{} takes advantage of a WiredTiger's storage optimizations and is designed to exploit the good read performance it offers.
We keep the reads as sequential as is possible and use WiredTiger's cursor operations to keep things sequential.\ugh{rephrase}


\textbf{In-Situ Graph Processing:} Databases reign supreme for the management of structured data. \PM{in-DB analytics, HTAP, graph DB. Also include: GraphGen. I am using the term in-situ wrong. Apparently it means analytics on raw data}
Inputs for graph processing systems can be considered a view on a relational database, and therefore, at first glance, databases appear to be good candidates for performing in-situ analytics.
Traditional databases such as Postgres and MySQL are typically not used for \emph{offline} analytics tasks because the SQL interface is not convenient to write iterative queries, because it forces the developer to reason about relational entity tables and not the graph itself.
This has led to development of several Domain-Specific Languages (DSL's) for writing graph queries.
Vertexica \cite{jindal2014vertexica}, GreenMarl\cite{10.1145/2248487.2151013}, and Grail \cite{fan2015case} are examples of systems that follow this paradigm.
They present a "think like a vertex" abstraction atop SQL to provide an intuitive programming style to the application developers and demonstrate that many optimizations can be employed to attain performance competitive with specialized graph processing systems.
Other systems approach the flip side of this problem - providing a SQL-like interface to perform structural queries on graphs that have been stored in a RDBMS.
Zhao et.al.\cite{10.1145/3035918.3035943} create special SQL-level joins that provide efficient solutions for supporting graph processing.
PGQL \cite{pgql2016} provides a SQL-like interface to perform structural queries on graphs that have been stored in a RDBMS.
These SQL/RDBMS systems greatly benefit from the many decades of research and experience in building and tuning query optimizers.
Sevenich et.al. design a graph database\cite{10.14778/3007263.3007265} that uses PGX.D \cite{hong2015pgx} to bulk load data from a RDBMS into memory before running graph analytic algorithms that are expressed in GreenMarl \cite{10.1145/2248487.2151013}.
This system is representative of a "tight integration of a graph storage and a graph analytics engine" and is a database-backed version of in-memory graph analytic system.

\textbf{Graph Native Databases}
Graph databases, unlike graph processing systems, focus on real-time querying and manipulation of entities, relationships, and the graph structure.
Unlike RDBMSes, graph databases treat relationships as first class citizens, use graph-native storage, and employ index-free adjacencey lists to enable efficient traversals.
The most widely used databases are Neo4j\cite{neo4j}, ArangoDB\cite{arangodb}, JanusGraph\cite{janusgraph}, Amazon Neptune \cite{vogt_2000}, and TigerGraph\cite{deutsch2019tigergraph}.
The work by McColl et. al. \cite{McColl2014} presents a thorough performance comparison of different graph databases when running benchmarks on small datasets (The largest dataset has 135M vertices).
Graph databases are designed to run online queries, the kind that have been standardized by the Linked Data Benchmark Council's (LDBC) Social Network Benchmark group \cite{ldbcsite}.
A recent study by Pacaci et. al. \cite{specialized_gdb_waterloo} concludes that there is much performance variation between commercial graph databases and that Postgres often outperforms them for real-time graph queries.

\PM{optional: graph representations, HTAP systems}